% sections/related.tex -- Owner: W1. Every \cite key is defined in ../lit/refs.bib,
% except the w2: keys, which live in extra.bib (used here to avoid duplicate
% bibliography entries for works W2 already cites). Evidence for each claim:
% ../lit/killreads.md (full-text reads, with page numbers) and ../lit/extraction.csv.
\section{Related Work}\label{sec:related}

\paragraph{Formal analyses of sender-constrained protocols.}
Hosseyni, K\"usters and W\"urtele prove the FAPI~2.0 family secure in the Web
Infrastructure Model, including DPoP, mTLS-bound tokens and message
signing~\cite{hosseyni2022tr,hosseyni2024csf,hosseyni2024tops}; earlier work
covered OAuth~2.0 and FAPI~1.0~\cite{fett2016oauth,fett2019fapi}. In these
models the resource server is a single atomic process. It builds the URI it
compares with \htu from the host and path of the message it serves, and it
selects the resource from the same path. The verified view therefore \emph{is}
the executed request by construction. No intermediary, URI normalisation or
method override appears, and the security property is confidentiality of the
resource rather than integrity of the executed operation. The FAPI~2.0
attacker model leaks resource requests only \emph{after} the resource server
has processed them~\cite{fapiam2025}. Our model keeps their checks and
splits the resource server into a verifier and an executor. Binding soundness
is the property their proofs assume implicitly.

\paragraph{HTTP interpretation differentials.}
A long line of work finds disagreements between HTTP components:
\begin{itemize}[leftmargin=*,nosep]
  \item request smuggling and desync~\cite{jabiyev2021treqs,jabiyev2024gudifu,mu2025hdhunter,kallus2024garden,shen2022hdiff};
  \item HTTP/2-to-HTTP/1 conversion~\cite{jabiyev2022frameshifter,gavazzi2025h2fuzz};
  \item Host ambiguity~\cite{chen2016host};
  \item URL parsing~\cite{reynolds2022equivocal};
  \item cache path confusion and poisoning~\cite{mirheidari2020cached,nguyen2019cpdos}.
\end{itemize}
WAFFLED is closest in structure: a WAF allows a request whose body the
framework then parses and executes differently~\cite{akhavani2025waffled}.
These works supply many of our \GB primitives, but their oracles are
framing, caching or content-filter outcomes. Ours is a cryptographic proof
that is accepted while a different operation executes. Two of our three
classes, \GA and \GC, need no parsing disagreement at all.

\paragraph{Defences against differentials.}
HTTP Request Synchronization makes each hop's interpretation of the URL, the
Host and the framing consistent along a proxy chain~\cite{topcuoglu2025sync}.
It does not synchronise the effective method, body-level dispatch or trust
headers, and it is not tied to the client's proof. Under \GA every hop
already agrees, so there is nothing to synchronise. Its artefact was not
publicly available when we checked. Other defences do not bind the client's
proof to the executed operation either:
\begin{itemize}[leftmargin=*,nosep]
  \item header whitelisting~\cite{buttner2021whitelisting};
  \item cryptographic hop binding~\cite{nygren2025binding};
  \item cache-aware message authentication~\cite{nguyen2020crehma};
  \item HTTP Message Signatures, whose coverage is a deployment choice~\cite{rfc9421}.
\end{itemize}
Our defences D-A and D-B (\S\ref{sec:defence}) bind exactly that.
XML Signature Wrapping is the classic precedent of ``verified element $\neq$
processed element''~\cite{somorovsky2012saml}. We show the HTTP analogue for
proof-of-possession.

\paragraph{Execution semantics and placement.}
Golinelli et al.\ measure HTTP method interchangeability and find CSRF
defences that are scoped to one method~\cite{golinelli2026post}. We study
method \emph{override} and query or body dispatch under a proof that binds
\htm, across gateway placements. Single-product advisories show that an edge
check can see a different path or method than the backend executes:
\begin{itemize}[leftmargin=*,nosep]
  \item path views~\cite{traefik2026cve40912,istio2021cve31920,opaenvoy2026cve26205,envoy2019ghsa};
  \item method override, which ``allows bypassing method-based access controls enforced by upstream proxies''~\cite{grpcgw2026cve37236}.
\end{itemize}
We systematise this class and show that a DPoP check inherits it.
Mancinas Fontalvo reports silent binding failures in sender-constrained
deployments~\cite{mancinas2026}: replay through bounded \jti caches, and lossy
client-certificate propagation for mTLS. That author argues DPoP is immune to
the latter because its binding lives at the application layer. By
construction, however, an application-layer binding is still relative to the
verifier's view of the request (\S\ref{sec:model:classes}).
Conceptual work also argues that accepted tokens can be
insufficient~\cite{gallo2026continuity,gallo2026relational,kim2026tbsg}.

\paragraph{Testing and measurement of OAuth deployments.}
Prior tools and studies examine different layers:
\begin{itemize}[leftmargin=*,nosep]
  \item OAuch tests authorization-server compliance~\cite{philippaerts2022oauch};
  \item JWTeemo tests JOSE libraries~\cite{yang2026jwteemo};
  \item recent studies measure MCP server authentication~\cite{zhou2026mcp,padilla2026mcp}.
\end{itemize}
The MCP conformance suite adds negative DPoP tests at a single
endpoint~\cite{mcpconf369}. None of them checks whether a proof minted for
one operation can drive another.

\paragraph{Standards discussion.}
RFC~9449 binds only the method and the query-less URI~\cite{rfc9449}.
Recent IETF OAuth WG threads propose relaxing or dropping \htu and \htm
because intermediaries rewrite URIs, or forwarding proofs to downstream
servers~\cite{allam2026thread,madden2025thread}. The MCP DPoP profile does not
require \jti tracking~\cite{sep1932}. FAPI~2.0 and the OAuth security BCP
mandate sender-constraining but say nothing about where it is
verified~\cite{fapisp2025,rfc9700}. Binding soundness gives these
discussions a criterion for what a relaxed or forwarded check must preserve.
